\documentclass{article}
% Language setting
% Replace `english' with e.g. `spanish' to change the document language
\usepackage[english]{babel}

\usepackage[letterpaper,% Set the height and width of the paper
% includehead,
nomarginpar,% We don't want any margin paragraphs
textwidth=15cm,% Set \textwidth to 15cm
headheight=0mm,% Set \headheight to 10mm
headsep=0mm,
]{geometry}

\usepackage{fancyhdr}
\fancyhead{} % clear all header fields
\fancyfoot{} % clear all footer fields
\renewcommand{\headrulewidth}{0pt}  % no line in header area
\renewcommand{\footrulewidth}{1pt}  % no line in header area
\pagestyle{fancy}
\fancyfoot[L]{\today}

% Useful packages
\usepackage{amsmath,amssymb}
\usepackage{graphicx}
\usepackage{xcolor}
\usepackage[colorlinks=true, allcolors=blue]{hyperref}

\usepackage{titlesec}
% \renewcommand{\contentsname}{}
\titleformat*{\section}{\bfseries}
\font\titlefont=cmr12 at 15pt
\font\authorfont=cmr12 at 10pt
\title{{\titlefont Errata: Practical Uncertainty in Neural Networks}}
\author{{\authorfont Ethan Goan}}
\date{} % no date

\DeclareMathOperator{\erf}{erf}

\begin{document}
\maketitle
\thispagestyle{fancy} % make footer appear

\textbf{p37, Eqn 3.30} for leap frog integration, the step should be $\epsilon/2$. I am not happy with the notation there. Instead use,

\begin{align*}
  v_i(t + \frac{\epsilon}{2}) = v_i(t) + \frac{\epsilon}{2} \dfrac{dv_{i}(t)}{dt},
  \\
  w_i(t + \epsilon) = w_i(t) + \epsilon \frac{dw_{i}(t + \epsilon /2)}{dt},
  \\
  v_i(t + \epsilon) = v_i(t +\frac{\epsilon}{2}) + \frac{\epsilon}{2} \dfrac{dv_{i}(t + \epsilon)}{dt}.
\end{align*}


\textbf{p51} For the $\alpha$ divergence on page 75, the limits to relate to forward/reverse KL divergence and Hellinger distance were not correct. It should be, $KL(q|p): \alpha \rightarrow 0, KL(p|q): \alpha \rightarrow 1,$ Hellinger: $ \alpha \rightarrow 0.5$.

\textbf{p51 Eq. 3.58} missing factor of $1/2$


\textbf{p107, Eq 4.7:}  $ p(\mathbf{w} | \mathbf{X}, \mathbf{Y} , \theta ) \propto \Big[\prod_{i=0}^{N-1}\mathcal{N}\Big({\color{red}\mathbf{y}}|f(\mathbf{x}_{i};\theta) \mathbf{w}, \sigma^{2}\Big)\Big] p(\mathbf{w}).$

\textbf{Eq 4.11 and 4.15:} are missing $\Phi$, they should read

$\sigma^{2}_{N}(\hat{\mathbf{x}}) =
\sigma^{2} + \Phi(\hat{\mathbf{x}})^{T}\Sigma_{\pi} \Phi(\hat{\mathbf{x}})$ and
$\sigma^{2}_{N}(\hat{\mathbf{x}}) \approx \sigma^{2}
\Phi(\hat{\mathbf{x}})^{T}\Sigma_{\text{diag}} \Phi(\hat{\mathbf{x}})$

\textbf{p111, Eq 4.22:} The approximation from ref [34] in the paper provides the Taylor Series approximation of
$ \mathcal{N}\Big(\frac{\mu_{j}}{\mu_{d}},\frac{\mu_{j}^{2}}{\mu_{d}^{2}}\big(\sigma^{2}_{j} / \mu_{j}^{2} +  \sigma^{2}_{d} / \mu_{d}^{2} \big)\Big)$. The issue with this is the square of the $\mu_{d}$ terms, as this is at risk of numerical overflow in half floating-point precision. Instead we use in $\sigma^{2} \approx \mu_{j} / \mu_{d}\sqrt{\sigma^{2}_{j} + \sigma^{2}_{d}}$, which emperically follwed a similar trend and was less of a risk for numerical overflow (though not immune).

\textbf{p135, Eq. 5.13}, the integral sign should be there in the denominator.


\textbf{p143}, Eqn 5.43, and \textbf{Pg. 146} Eqn 5.46, Typo in missing square term.
  \begin{align*}
    \frac{\sigma^{2}}{2} &\int_{-\infty} ^{\infty}v^{2} \erf\Big(\frac{(\sigma v + \mu)}{\sqrt{2}}\Big)\phi(v) dv \nonumber
    \\ = &\frac{\sigma^{2}}{4\sqrt{2\pi}} \Big[ \frac{\sqrt{\pi}}{\alpha_{1} \sqrt{\alpha_{1}}}
    \erf\Big({\frac{\sqrt{\alpha_{1}}b_{1}}{\sqrt{\alpha_{1} + a_{1}^{2}}}}\Big)
    -  \frac{1}{\alpha_{1}\sqrt{\alpha_{1} + a_{1}^{2}}} \frac{2a_{1}^{2}b_{1}}{\alpha_{1} + a_{1}^{2}}
    \exp{\big(\frac{-\alpha_{1} b_{1}^{\textcolor{red}{2}}}{\alpha_{1} + a_{1}^{2}}\big)} \Big]
  \end{align*}

  This same typo carried over to integral property from [240, EQ. 2.7.3.9]
  \begin{align*}
    \int_{_{-\infty}}^{\infty} z^{2}\exp{(\alpha^{2})} \erf(\alpha_{1}z + \beta_{1})
    &=
    \frac{\sqrt{\pi}}{2\alpha \sqrt{\alpha}}
      \erf\Big({\frac{\sqrt{\alpha}\beta_{1}}{\sqrt{\alpha + \alpha_{1}^{2}}}}\Big)
      \nonumber\\&
    -  \frac{1}{\sqrt{\alpha + \alpha_{1}^{2}}} \frac{2\alpha_{1}^{2} \beta_{1}}{\alpha + \alpha_{1}^{2}}
    \exp{\frac{-\alpha \beta_{1}^{\textcolor{red}{2}}}{\alpha + \alpha^{2}}}
  \end{align*}


  And to the final result on \textbf{Pg. 147}, Eqn 5.61 for the approximation to the variance of the Gelu function,

  \begin{align*}
    \sigma_{g}^2 \approx & \frac{\sigma^{2} + \mu^{2}}{2} \nonumber
                 \\
    % second first term
              &+
                \frac{\sigma^{2}}{4\sqrt{2\pi}} \Big[ \frac{\sqrt{\pi}}{\alpha_{1} \sqrt{\alpha_{1}}}
                \erf\Big({\frac{\sqrt{\alpha_{1}}b_{1}}{\sqrt{\alpha_{1} + a_{1}^{2}}}}\Big)
                -  \frac{1}{\alpha_{1}\sqrt{\alpha_{1} + a_{1}^{2}}} \frac{2a_{1}^{2}b_{1}}{\alpha_{1} + a_{1}^{2}}
                \exp{\big(\frac{-\alpha_{1} b_{1}^{\textcolor{red}{2}}}{\alpha_{1} + a_{1}^{2}}\big)} \Big] \nonumber
    \\% second second term
              & +
                \frac{\mu \sigma a_{2}}{\alpha_{2}\sqrt{2\pi(a_{2}^{2} + \alpha_{2})}}
                \exp\Big( - \frac{\alpha_{2} b_{2}^{2}}{a_{2}^{2} + \alpha_{2}} \Big) \nonumber
    \\% second third term
              & +
                \frac{\mu^{2}}{2\sqrt{2\alpha_{3}}} \erf\Big(
                \frac{2\alpha_{3}b_{3}}{2\sqrt{\alpha_{3}^{2} + \alpha_{3}a_{3}^{2}}} \Big) \nonumber
    \\% third term
              & - \frac{(2\alpha_{4}^{2} + b_{4}^{2})}{16\sigma\alpha_{4}^{5}\sqrt{2}}
                \exp \Big(
                \frac{b_{4}^{2}}{4\alpha_{4}^{2}} + \gamma \Big)
                - \mu_{g}^{2}
  \end{align*}
  where, $\alpha_{1}, \alpha_{2}, \alpha_{3} = \frac{1}{2}$, $b_{1},
    b_{2}, b_{3} = \frac{\mu}{\sqrt{2}}$, $a_1, a_{2},
    a_{3} = \frac{\sigma}{\sqrt{2}}$, $\alpha_{4} =
    \sqrt{\frac{1}{2\sigma^{2}} + \frac{\eta}{2}}$, $b_{4} =
    \frac{\mu}{\sigma^{2}}$ and $\gamma = -\frac{\mu^{2}}{2\sigma^{2}}$, $\eta = 1.23907$.

\textbf{pg. 150, Eqn 5.68}, the last term are wrong way around. Should read,
 \begin{equation*}
    \Omega = \Omega_{X}\Omega_{Y} + \Omega_{X}\mu^{2}_{Y} + \textcolor{red}{\mu^{2}_{X}\Omega_{Y}}.
  \end{equation*}

\end{document}
